\documentclass[11pt,a4paper]{article}
\usepackage[utf8]{inputenc}
\usepackage[T1]{fontenc}
\usepackage{lmodern}
\usepackage[margin=0.85in, headheight=16pt, top=1.1in, bottom=1.1in]{geometry}
\usepackage{amsmath,amssymb,amsfonts}
\usepackage{booktabs}
\usepackage{tabularx}
\usepackage{array}
\usepackage{hyperref}
\usepackage{xcolor}
\usepackage{listings}
\usepackage{fancyhdr}
\usepackage{titlesec}
\usepackage{microtype}
\usepackage{parskip}
\usepackage{caption}

% --- Professional Color Palette ---
\definecolor{primaryblue}{RGB}{18, 52, 86}
\definecolor{accentblue}{RGB}{0, 114, 206}
\definecolor{darkgray}{RGB}{35, 35, 35}
\definecolor{lightgraybg}{RGB}{248, 249, 250}
\definecolor{codebg}{RGB}{245, 247, 250}
\definecolor{keywordcolor}{RGB}{0, 102, 153}
\definecolor{stringcolor}{RGB}{163, 21, 21}
\definecolor{commentcolor}{RGB}{0, 128, 0}

% --- Tabularx Column Types ---
\newcolumntype{Y}{>{\raggedright\arraybackslash}X}

% --- Hyperlink Setup ---
\hypersetup{
    colorlinks=true,
    linkcolor=primaryblue,
    urlcolor=accentblue,
    citecolor=primaryblue,
    pdfauthor={Lead Systems Engineer},
    pdftitle={System Engineering Specification}
}

% --- Code Listing Setup ---
\lstset{
    backgroundcolor=\color{codebg},
    basicstyle=\ttfamily\scriptsize\color{darkgray},
    breaklines=true,
    breakatwhitespace=true,
    columns=fullflexible,
    keepspaces=true,
    numbers=left,
    numberstyle=\tiny\color{gray!70},
    numbersep=5pt,
    frame=single,
    framerule=0.5pt,
    rulecolor=\color{primaryblue!30},
    commentstyle=\color{commentcolor}\itshape,
    keywordstyle=\color{keywordcolor}\bfseries,
    stringstyle=\color{stringcolor},
    xleftmargin=12pt,
    framexleftmargin=12pt,
    tabsize=4,
    showstringspaces=false
}

% --- Header & Footer (Collision-Proof) ---
\pagestyle{fancy}
\fancyhf{}
\setlength{\headheight}{16pt}
\fancyhead[L]{\small\textcolor{primaryblue}{\textbf{Project Technical Specification}}}
\fancyhead[R]{\small\textcolor{gray}{Engineering Dept.}}
\fancyfoot[C]{\small\thepage}
\fancyfoot[R]{\tiny\textcolor{gray!70}{DOC-SPEC-001}}
\renewcommand{\headrulewidth}{0.4pt}

% --- Professional Headings ---
\titleformat{\section}{\large\bfseries\color{primaryblue}}{\thesection}{0.8em}{}[\vspace{-2pt}\titlerule]
\titleformat{\subsection}{\normalsize\bfseries\color{primaryblue!90}}{\thesubsection}{0.8em}{}
\titleformat{\subsubsection}{\small\bfseries\color{primaryblue!75}}{\thesubsubsection}{0.8em}{}
\titlespacing*{\section}{0pt}{10pt}{4pt}
\titlespacing*{\subsection}{0pt}{8pt}{3pt}

\title{
    \vspace{-1.2cm}
    \textbf{\huge \color{primaryblue} Autonomous Robotics Platform} \\
    \vspace{0.2cm}
    \Large \textbf{System Engineering \& Architecture Specification} \\
    \vspace{0.15cm}
    \normalsize Document ID: \texttt{DOC-SPEC-001} \quad | \quad Revision 1.0.0
}
\author{\textbf{Lead Systems Engineer} \quad \& \quad \textbf{Robotics Engineering Team}}
\date{\today}

\begin{document}

\maketitle
\thispagestyle{fancy}

\begin{abstract}
\noindent This document establishes the official technical baseline, mathematical derivations, communication interfaces, and verification requirements for the robotic platform. The architecture ensures high-throughput telemetry ingestion, deterministic state estimation, and reliable integration across embedded compute platforms.
\end{abstract}

\vspace{-0.2cm}
\tableofcontents
\vspace{0.3cm}

\section{System Architecture}
The system architecture separates high-frequency real-time sensor polling from compute-intensive localization and computer vision workloads.

\begin{center}
\begin{lstlisting}[language={}, numbers=none, basicstyle=\ttfamily\tiny\color{darkgray}]
 [Sensor Suite] ---> [Real-time Microcontroller] ---> (Serial/CAN/RS485) ---> [ROS 2 Compute Host]
                                                                                     |
                                                                                     v
                                                                             [Navigation Core]
\end{lstlisting}
\vspace{-4pt}
\captionof{figure}{System Data Flow}
\end{center}

\section{Kinematics \& Mathematical Formulation}
The differential kinematics for linear velocity $v$ and angular velocity $\omega$ are defined by track velocities $v_L$ and $v_R$ separated by baseline gauge $W$:
\begin{equation}
v = \frac{v_R + v_L}{2}, \quad \omega = \frac{v_R - v_L}{W}
\end{equation}

For discrete integration over timestep $\Delta t$, the updated 2D pose $(x_{k+1}, y_{k+1}, \theta_{k+1})$ is calculated as:
\begin{equation}
\theta_{k+1} = \theta_k + \omega \Delta t, \quad x_{k+1} = x_k + v \Delta t \cos\theta_{k+1}, \quad y_{k+1} = y_k + v \Delta t \sin\theta_{k+1}
\end{equation}

\section{Interface Control Document (ICD)}
\begin{center}
\footnotesize
\renewcommand{\arraystretch}{1.15}
\begin{tabularx}{\textwidth}{@{} l l c Y @{}}
\toprule
\textbf{Topic / Interface} & \textbf{Message Type} & \textbf{Rate} & \textbf{Description} \\
\midrule
\texttt{/odom} & \texttt{nav\_msgs/Odometry} & 100\,Hz & Position $(X,Y,Z)$, quaternion, linear and angular velocities \\
\texttt{/robot/pose} & \texttt{geometry\_msgs/PoseStamped} & 100\,Hz & 2D/3D stamped pose in \texttt{odom} coordinate frame \\
\texttt{/robot/sensors} & \texttt{sensor\_msgs/Imu} & 100\,Hz & Calibrated 9-DoF IMU acceleration and gyro rates \\
\texttt{/tf} & \texttt{tf2\_msgs/TFMessage} & 100\,Hz & Dynamic transformation tree (\texttt{odom} $\to$ \texttt{base\_link}) \\
\bottomrule
\end{tabularx}
\vspace{-4pt}
\captionof{table}{System Interface Specification}
\end{center}

\section{Verification \& Test Summary}
All software components are subjected to continuous integration and automated testing:
\begin{itemize}
    \item \textbf{Unit Tests}: $100\%$ pass rate across kinematic integration suites.
    \item \textbf{Timing Jitter}: Telemetry packet arrival latency guaranteed $< 1.5\,\text{ms}$.
\end{itemize}

\end{document}
